\textbf{Phần 1. Trắc nghiệm lựa chọn}
\setcounterpageref{ex}{0}

\begin{ex} %[2D7V34]
Tính giá trị biểu thức: $46 + 28 - 35$.
\choice
{$37$}
{$38$}
{$40$}
{\True $39$}
\loigiai{
$46 + 28 = 74$; $74 - 35 = 39$.
\textbf{Đáp án của câu hỏi là D}
}
\end{ex}

\begin{ex} %[2D7V34]
Tìm $x$ biết: $x - 29 = 45$.
\choice
{$72$}
{\True $74$}
{$73$}
{$75$}
\loigiai{
$x = 45 + 29 = 74$.
\textbf{Đáp án của câu hỏi là B}
}
\end{ex}

\begin{ex} %[2D7V34]
Một đàn gà có $52$ con, mẹ bán đi $18$ con rồi mua thêm $15$ con nữa. Hỏi đàn gà có bao nhiêu con?
\choice
{$48$ con gà}
{$50$ con gà}
{\True $49$ con gà}
{$47$ con gà}
\loigiai{
$52 - 18 + 15 = 34 + 15 = 49$ (con gà).
\textbf{Đáp án của câu hỏi là C}
}
\end{ex}

\begin{ex} %[2D7V34]
Tính nhanh: $36 + 18 + 24 + 12$.
\choice
{\True $90$}
{$85$}
{$80$}
{$95$}
\loigiai{
$(36 + 24) + (18 + 12) = 60 + 30 = 90$.
\textbf{Đáp án của câu hỏi là A}
}
\end{ex}

\begin{ex} %[2D7V34]
Một sợi dây dài $80\text{ cm}$. Người ta cắt lấy $2$ đoạn, mỗi đoạn dài $25\text{ cm}$. Đoạn dây còn lại dài bao nhiêu xăng-ti-mét?
\choice
{$25\text{ cm}$}
{$35\text{ cm}$}
{$20\text{ cm}$}
{\True $30\text{ cm}$}
\loigiai{
Hai đoạn dài: $25 + 25 = 50\text{ (cm)}$. Còn lại: $80 - 50 = 30\text{ (cm)}$.
\textbf{Đáp án của câu hỏi là D}
}
\end{ex}

\begin{ex} %[2D7V34]
Hai số có tổng bằng $68$, nếu số hạng thứ nhất tăng $7$ đơn vị và số hạng thứ hai giảm $5$ đơn vị thì tổng mới là
\choice
{$69$}
{\True $70$}
{$71$}
{$66$}
\loigiai{
Tổng mới: $68 + 7 - 5 = 70$.
\textbf{Đáp án của câu hỏi là B}
}
\end{ex}

\begin{ex} %[2D7V34]
Điền số thích hợp: $75 - \dots = 38 + 17$.
\choice
{$19$}
{$21$}
{\True $20$}
{$18$}
\loigiai{
Vế phải: $38 + 17 = 55$. Ô trống $= 75 - 55 = 20$.
\textbf{Đáp án của câu hỏi là C}
}
\end{ex}

\begin{ex} %[2D7V34]
Lớp 2A có $35$ học sinh, lớp 2B có $32$ học sinh. Cả hai lớp có bao nhiêu học sinh?
\choice
{\True $67$ học sinh}
{$65$ học sinh}
{$66$ học sinh}
{$68$ học sinh}
\loigiai{
$35 + 32 = 67$ (học sinh).
\textbf{Đáp án của câu hỏi là A}
}
\end{ex}

\begin{ex} %[2D7V34]
Trong một phép trừ, số bị trừ là $82$, hiệu là $39$. Số trừ là
\choice
{$41$}
{$42$}
{$44$}
{\True $43$}
\loigiai{
Số trừ: $82 - 39 = 43$.
\textbf{Đáp án của câu hỏi là D}
}
\end{ex}

\begin{ex} %[2D7V34]
Bác Ba thu hoạch được $72\text{ kg}$ cam. Bác bán được $48\text{ kg}$. Bác còn lại bao nhiêu ki-lô-gam cam?
\choice
{$22\text{ kg}$}
{\True $24\text{ kg}$}
{$23\text{ kg}$}
{$25\text{ kg}$}
\loigiai{
$72 - 48 = 24\text{ (kg)}$.
\textbf{Đáp án của câu hỏi là B}
}
\end{ex}

\begin{ex} %[2D7V34]
Tìm một số biết rằng lấy số đó cộng với $34$ thì bằng $91 - 18$.
\choice
{$38$}
{$40$}
{\True $39$}
{$37$}
\loigiai{
Vế phải: $91 - 18 = 73$. Số đó là: $73 - 34 = 39$.
\textbf{Đáp án của câu hỏi là C}
}
\end{ex}

\begin{ex} %[2D7V34]
Điền dấu thích hợp: $56 + 27 \dots 92 - 15$.
\choice
{\True $>$}
{$<$}
{$=$}
{Không so sánh được}
\loigiai{
$56 + 27 = 83$; $92 - 15 = 77$. Vì $83 > 77$ nên điền dấu $>$. 
\textbf{Đáp án của câu hỏi là A}
}
\end{ex}

\textbf{Phần 2. Trắc nghiệm đúng sai}
\setcounterpageref{ex}{0}

\begin{ex} %[2D7V34]
Xét tính đúng sai của các biểu thức tính toán:
\choiceTF
{\True Biểu thức $38 + 47 - 25$ có kết quả là $60$.}
{Hiệu của $72 - 38$ bằng hiệu của $82 - 38$.}
{\True Trong phép cộng, tổng luôn lớn hơn hoặc bằng mỗi số hạng.}
{\True Phép tính $65 - 28$ có kết quả là số lẻ.}
\loigiai{
\begin{itemize}
    \item Ý a) Đúng: $38 + 47 = 85$; $85 - 25 = 60$.
    \item Ý b) Sai: Hai số bị trừ khác nhau thì hiệu phải khác nhau.
    \item Ý c) Đúng: Đúng với các số tự nhiên.
    \item Ý d) Đúng: $65 - 28 = 37$, là số lẻ. Đúng.
\end{itemize}
\begin{itemize}
    \item a) \textbf{Đúng}.
    \item b) \textbf{Sai}.
    \item c) \textbf{Đúng}.
    \item d) \textbf{Đúng}.
\end{itemize}
}
\end{ex}

\begin{ex} %[2D7V34]
Một cửa hàng có 95 mét vải, ngày thứ nhất bán 38 mét, ngày thứ hai bán 29 mét:
\choiceTF
{\True Tổng số mét vải bán trong hai ngày là 67 mét.}
{\True Cửa hàng còn lại 28 mét vải.}
{Ngày thứ hai bán nhiều hơn ngày thứ nhất.}
{Số mét vải còn lại nhiều hơn số mét vải bán ngày thứ hai.}
\loigiai{
\begin{itemize}
    \item Ý a) Đúng: $38 + 29 = 67\text{ m}$.
    \item Ý b) Đúng: $95 - 67 = 28\text{ m}$.
    \item Ý c) Sai: Ngày thứ hai (29 m) ít hơn ngày thứ nhất (38 m).
    \item Ý d) Sai: Còn 28 m, bán 29 m nên còn ít hơn.
\end{itemize}
\begin{itemize}
    \item a) \textbf{Đúng}.
    \item b) \textbf{Đúng}.
    \item c) \textbf{Sai}.
    \item d) \textbf{Sai}.
\end{itemize}
}
\end{ex}

\textbf{Phần 3. Trả lời ngắn}
\setcounterpageref{ex}{0}

\begin{ex} %[2D7V34]
Tính giá trị: $58 + 37 - 49$.
\shortans{46}
\loigiai{
\textbf{Đáp số:} $46$\\
$58 + 37 = 95$; $95 - 49 = 46$.
}
\end{ex}

\begin{ex} %[2D7V34]
Tìm x biết: $91 - x = 36 + 18$.
\shortans{37}
\loigiai{
\textbf{Đáp số:} $37$\\
$36 + 18 = 54 \Rightarrow x = 91 - 54 = 37$.
}
\end{ex}

\begin{ex} %[2D7V34]
Một đàn vịt có 65 con, người ta bán đi 28 con. Hỏi đàn vịt còn lại bao nhiêu con?
\shortans{37}
\loigiai{
\textbf{Đáp số:} $37$\\
$65 - 28 = 37$ con.
}
\end{ex}

\begin{ex} %[2D7V34]
Tổng của hai số là 82. Nếu bớt ở số thứ nhất đi 15 đơn vị thì tổng mới bằng bao nhiêu?
\shortans{67}
\loigiai{
\textbf{Đáp số:} $67$\\
$82 - 15 = 67$.
}
\end{ex}

\begin{ex} %[2D7V34]
Tìm một số biết rằng lấy số đó trừ đi 38 rồi cộng thêm 19 thì được 52.
\shortans{71}
\loigiai{
\textbf{Đáp số:} $71$\\
Số đó: $52 - 19 + 38 = 71$.
}
\end{ex}

\begin{ex} %[2D7V34]
Hai bạn có tất cả 54 viên bi. Nếu bạn An cho bạn Bình 6 viên bi thì hai bạn có số bi bằng nhau. Hỏi lúc đầu bạn An có bao nhiêu viên bi?
\shortans{33}
\loigiai{
\textbf{Đáp số:} $33$\\
Khi bằng nhau mỗi bạn có: $54 : 2 = 27$ viên. Ban đầu An có: $27 + 6 = 33$ viên.
}
\end{ex}
